% Generated draft: requires mathematical and source review.
\documentclass[11pt]{article}
\usepackage[margin=25mm]{geometry}
\usepackage{fontspec}
\setmainfont{Latin Modern Roman}
\usepackage{amsmath,amssymb,mathtools,unicode-math}
\setmathfont{Latin Modern Math}
\usepackage{enumitem,needspace,xurl,hyperref}
\setlength{\parindent}{0pt}
\setlength{\parskip}{5pt}
\setlength{\emergencystretch}{3em}
\newcommand{\N}{\mathbb{N}}
\newcommand{\Z}{\mathbb{Z}}
\newcommand{\Q}{\mathbb{Q}}
\newcommand{\R}{\mathbb{R}}
\newcommand{\C}{\mathbb{C}}
\newcommand{\F}{\mathbb{F}}
\newcommand{\A}{\mathbb{A}}
\newcommand{\PP}{\mathbb{P}}
\newcommand{\E}{\mathbb{E}}
\newcommand{\eps}{\varepsilon}
\newcommand{\dd}{\,\mathrm{d}}
\begin{document}
\section*{UnsolvedMath ID: 10000022}
\textbf{Source problem code:} catalogue entry
\subsection{Definitions and mathematical statement}
Let $f,g:\mathbb N\to\mathbb Z$ and retain the literal set in the source:
\[
S=\{(n,f(n),g(n)):n\in\mathbb N\}\subset\mathbb Z^3.
\]
A usual exponential-intersection-tails property asks for a probability law $\mu$ on infinite self-avoiding nearest-neighbor paths from a fixed vertex such that, for independent paths $\gamma,\gamma'$ with law $\mu$, there are constants $C<\infty$ and $a>0$ satisfying
\[
\Pr\bigl(|\operatorname{range}(\gamma)\cap\operatorname{range}(\gamma')|\geq m\bigr)
\leq Ce^{-am}\qquad(m\geq1).
\]
The constants do not depend on $m$. The source does not specify this path/root convention or whether intersections count vertices or edges, so this is explicit standard context.

Characterize the functions $f,g$ for which the intended path law exists. There is a serious geometric ambiguity: the displayed $S$ has exactly one vertex at each first-coordinate level, whereas the title describes three-dimensional slabs. In its induced nearest-neighbor graph, two distinct vertices can be adjacent only at consecutive levels where both $f$ and $g$ are unchanged. Thus this literal set is not a thick slab and cannot silently be treated as one. The intended subset/adjacency, path class, and starting-point convention must be checked before deriving a slab-growth criterion.
\subsection{Short English statement}
Determine when the source's subset supports a law of infinite paths with exponentially decaying pairwise intersection tails. The displayed one-point-per-level set conflicts with the slab description, and path/adjacency conventions are missing.
\subsection{Sources}
\begin{itemize}
\item Itai Benjamini, \emph{Coarse Geometry and Randomness}, Saint-Flour notes, Open Problem 4.30, PDF page 37, \url{https://www.wisdom.weizmann.ac.il/~itai/stflouraug24.pdf#page=37}; preserved copy \url{https://arquivo.pt/noFrame/replay/20201231041548id_/http://www.wisdom.weizmann.ac.il/~itai/stflouraug24.pdf}; accessed 2026-07-29. The supplied provenance reports extracted text checked against rendered pages.
\item Local registry \texttt{data/unsolvedmath/problems.json}, permanent UnsolvedMath ID $10000022$, source code catalogue entry, \emph{Exponential intersection tails in three-dimensional slabs}. Dataset attribution: Ulam AI and the UnsolvedMath Contributors.
\end{itemize}
\end{document}
